% Corsin Week 5 -- full exercise sheet (complete; important problems 5.1, 5.2, 5.3 are also
% duplicated inline at the top of content/week-05.tex). NOT \input by main.tex -- reference
% source only.
\exercisesheet{5}

\textit{Statements quoted verbatim from \texttt{exercises/Ex5\_Analysis2\_eng.pdf} (assigned 16
March 2026, due 23 March 2026). Attribution: Prof. Joaquim Serra, D-MATH, ETH Z\"urich.}

\begin{notation}[Corsin's recommendations]
\begin{center}
\begin{tabular}{lll}
\textbf{Problem} & \textbf{Priority} & \textbf{Corsin's note} \\
\hline
5.1 & important & ``Recall the linearity of the differential and that $C^1 \Rightarrow$ \\
    & & differentiable'' \\
5.2 & important & ``Compare with the example from the script and recall that continuous \\
    & & functions on compact sets attain their extrema.'' \\
5.3 & important & ``There is no difference between $x \to 0$ and $|x| \to 0$.'' \\
5.4 & semi-important & ``Use the definition of directional derivative, \\
    & & $\partial_v f(x) = \left.\frac{\partial}{\partial t}\right\|_{t=0} f(x+tv)$. Solve a \\
    & & few of these, but not all.'' \\
5.5 & optional & --- \\
5.6 & optional & ``Interesting, but irrelevant for the course'' \\
5.7 & optional & ``Interesting for physicists, will be covered in classical mechanics.'' \\
\end{tabular}
\end{center}
\end{notation}

\begin{exercise}[5.1 --- Derivatives \textnormal{(important)}]
\label{ex:5.1}
\begin{enumerate}[label=\textbf{(\alph*)}]
  \item Consider $f(x,y) := xy^2e^{-x^2-y^2}$ for $(x,y) \in \mathbb{R}^2$. Find all the critical
    points of $f$, that is all the points where the gradient of $f$ vanishes. (The point will be
    given if and only if all the numerical values you find are correct\ldots\ so check your
    computations twice.)
  \item We want to find a function $g \in C^1(\mathbb{R}^2)$ with the following directional
    derivative:
    $$\partial_v g(x,y) = 2\cos(x^2y)xv_2 + \cos(x^2y)v_1^2y \quad \text{for all } (x,y) \text{ and } (v_1,v_2) \in \mathbb{R}^2.$$
    Give an explicit example of such a function $g$ or prove that a function with these
    properties cannot exist.
\end{enumerate}
\end{exercise}

\begin{remark}
Corsin's hint --- \textit{recall the linearity of the differential} --- is the whole of
\cref{ex:5.1}(b): a directional derivative must be \textbf{linear} in $v$, and $v_1^2$ is not.
\end{remark}

\begin{exercise}[5.2 --- Lagrange multipliers \textnormal{(important)}]
\label{ex:5.2}
Consider the ball $U := \{(x,y,z) : x^2+y^2+z^2 < 1\}$ and the function
$$f(x,y,z) := 1 - x^2 - y^2 + 4x.$$
\begin{enumerate}[label=\textbf{\arabic*.}]
  \item Motivate rigorously why $f$ attains its absolute maximum and minimum in $\overline{U}$.
  \item Find all the critical points of $f$ which lie in the interior of $U$. Say whether
    $\sup_U f$ (or $\inf_U f$) is attained at some point in $U$.
  \item Say whether $\max_{\partial U} f$ and $\min_{\partial U}$ exist.
  \item With the method of Lagrange multipliers, find the possible critical points for the
    constrained problem $\max\{f(x,y,z) : (x,y,z) \in \partial U\}$, and same for min.
  \item Among all the points you have found, clarify at which points the maximum/minimum of $f$
    is attained.
\end{enumerate}
\end{exercise}

\begin{exercise}[5.3 --- Multiple choice (Landau notation) \textnormal{(important)}]
\label{ex:5.3}
Mark all and only the statements which are true.
\begin{enumerate}[label=\textbf{(\alph*)}]
  \item As $|x| \to 0$, if $f(x) = x_2x_1^2 + O(x_1^2)$, then $f(x) = O(|x|^2)$.
  \item As $x \to 0$, if $f(x) = x_2x_1^2 + O(x_1^3)$, then $f(x) = O(x_1^3)$.
  \item As $|x| \to 0$, if $f(x) = x_2x_1 + O(|x|^2)$, then $f$ is differentiable at $x = 0$.
  \item As $x \to 0$, if $f(x) = x_2x_1 + O(|x|^2)$, then $f$ is twice differentiable around
    $x = 0$.
  \item As $|x| \to 0$, if $f(x) = x_2x_1 + O(|x|^3)$, then $f$ is twice differentiable around
    $x = 0$ and $\partial_{11}f(0) = 0$.
\end{enumerate}
\textit{Official hint:} revise \textit{Corollary 10.45}.
\end{exercise}

\begin{exercise}[5.4 --- Computation of derivatives \textnormal{(semi-important)}]
\label{ex:5.4}
For each of the following functions compute the directional derivative in the general direction
$v$ at a general point $x$, that is $\partial_v f(x)$, where applicable.
\begin{enumerate}[label=\textbf{\arabic*.}]
  \item $x_1/x_2$, $e^{-x_1/x_2}$ for $x \in \mathbb{R}^2\setminus\{x_2 = 0\}$. (Also: can they be
    extended continuously or differentiably across $\{x_2 = 0\}$?)
  \item $e^{x_1x_2}\sin(x_1+x_2^2)$, $\dfrac{x_1^2}{x_1^2+x_2^2}$,
    $\dfrac{x_1^2x_2}{x_1^2+x_2^2}$, $\dfrac{x_1^2x_2}{x_1^2+x_2^4}$ for
    $x \in \mathbb{R}^2\setminus\{0\}$. (Also: can they be extended continuously or
    differentiably across the origin?)
  \item $|x|$, $|x|^\alpha$, $x/|x|$, $g(|x|)$, $x\cdot e$, $g(x\cdot e)$, for
    $\alpha \in \mathbb{R}$, $e \in \mathbb{R}^n$ and $g \in C^1(\mathbb{R})$.
  \item $ax$, $x\transp$, $\Tr(x)$, $x^2$, $x^3$, $\Tr(x^2)$ for $x, a \in \mathbb{R}^{n\times n}$.
  \item \textbf{(*)} $x^{-1}$, $x^{-2}$ for $x \in \mathbb{R}^{n\times n}$ with $\det x \neq 0$.
\end{enumerate}
\textit{Official hints:} to prove that a function \textbf{is} differentiable at 0 you have two
tools --- the definition of differentiability, and the sufficient condition of
\textit{Theorem 10.16}. To prove that a function \textbf{is not} differentiable at 0 you have
two tools --- show that the necessary condition of \textit{Proposition 10.14} does not hold, or
show that for some $C^1$ curve $\gamma : (-\delta,\delta)\to\mathbb{R}^n$ the function
$(f\circ\gamma)(t)$ is \textbf{not} differentiable at $t = 0$ (simpler, as you study a function of
one variable). \textbf{5.4.3} --- these functions are $C^\infty$, so compute them on $x+tv$ and
find an expansion in powers of $t$; by \textit{Corollary 10.45} it must be the Taylor polynomial,
from which the derivatives can be read off. For example $(x+tv)^2 = x^2 + (xv+vx)t + O(t^2)$.
\end{exercise}

\begin{exercise}[5.5 --- Multiple choice (cost of a Taylor expansion) \textnormal{(optional)}]
\label{ex:5.5}
Assume $g \in C^\infty(\mathbb{R})$ and $f \in C^\infty(\mathbb{R}^n)$ is such that
$$f(x) = x_1 + x_2x_1 + O(|x|^4) \quad \text{as } x \to 0.$$
We want to compute the Taylor expansion of $(g\circ f)$ at $x = 0$ of the highest possible order
with the information we have. We can ask for the value $g^{(k)}(t)$ for arbitrary
$k \in \mathbb{N}$, $t \in \mathbb{R}$, but we have to pay 5 CHF each time.

For a general $g$, which is the degree of the best (i.e.\ highest-order) Taylor polynomial we can
compute? How expensive will it be to compute it?
\begin{enumerate}[label=\textbf{(\alph*)}]
  \item degree 3 and 20 CHF
  \item degree 2 and 15 CHF
  \item degree 2 and 10 CHF
  \item degree 3 and 15 CHF
\end{enumerate}
\end{exercise}

\begin{exercise}[5.6 --- Polynomials \textnormal{(optional --- ``interesting, but irrelevant for the course'')}]
\label{ex:5.6}
Let $P, Q \in \mathbb{R}[x_1,\dots,x_n]$ be polynomials of $n$ variables; assume that there are
positive numbers $M, \sigma$ such that
$$|P(x) - Q(x)| \leq M|x|^\sigma \quad \text{for all } |x| \leq 1.$$
Show that $P$ and $Q$ have the same coefficients of order smaller than $\sigma$. That is, if we
write $P(X) = \sum_{\alpha\in\mathbb{N}^n} a_\alpha X^\alpha$,
$Q(X) = \sum_{\alpha\in\mathbb{N}^n} b_\alpha X^\alpha$, then $a_\alpha = b_\alpha$ for all
$|\alpha| < \sigma$.
\end{exercise}

\begin{exercise}[5.7 --- Chain rule for a curve, or: a car on Earth \textnormal{(optional --- ``interesting for physicists'')}]
\label{ex:5.7}
Let $R > 0$ and $\omega \in \mathbb{R}$ be constants. Let
$\sigma : I \subset \mathbb{R}\to\mathbb{R}^2$, $\sigma(t) = (a(t), b(t))$, which we interpret as
a curve in a ``map'' representing the longitude and latitude of a car at time $t$.

Define the ``rotating'' spherical coordinates $F : \mathbb{R}^2\times\mathbb{R}\to\mathbb{R}^3$,
$$F((x,y),t) = \big(R\cos y\cos(\omega t + x),\ R\cos y\sin(\omega t + x),\ R\sin y\big),$$
representing what a point on the map looks like on Earth from a satellite (thus we consider the
Earth as a rotating sphere with speed given by $\omega$, and we forget about translation around
the Sun) at time $t$.

Consider $r : \mathbb{R}\to\mathbb{R}^3$, $r(t) = F(\sigma(t), t)$, which is the position of the
car at time $t$.
\begin{enumerate}[label=\textbf{\arabic*.}]
  \item Compute the acceleration $r''(t)$ using the chain rule.
  \item Assume that the acceleration is always proportional to the radial vector $r(t)$, i.e.\
    $r''(t) = \lambda(t)\,r(t)$ for some scalar function $\lambda(t)$. Show that $a$ and $b$
    satisfy
    $$b'' + \sin b\cos b\,(\omega+a')^2 = 0, \qquad \cos b\,a'' - 2\sin b\,(\omega+a')b' = 0.$$
    Equivalently, away from the poles ($\cos b \neq 0$),
    $$a'' - 2\tan b\,(\omega+a')\,b' = 0.$$
\end{enumerate}
\end{exercise}
